% GENERATED FILE. Edit canonical lessons, then run npm run generate:books.
% canonical-source-hash: fnv1a64:3be9f4ef28e63316
% canonical-lessons: SA-C457-smayate, SA-C457-upahasati, SA-C457-parihasah, SA-C457-attahasah, SA-C457-vinodah

\chapter{Laughter and Jokes}
\label{ch:sa-laughter-and-jokes}
\hlchaptermodality{457}

\begin{chapteropening}
\textbf{By the end of this chapter:} \emph{I can say what made people laugh.}

Complete the last lesson of chapter 457: I can say what made people laugh.
\end{chapteropening}

\section[\texorpdfstring{\sk{स्मयते} (smayate)}{स्मयते (smayate)}]{\sk{स्मयते} (smayate) — smiles}

\label{lesson:SA-C457-smayate}



\begin{quote}
Before the new one: say the Sanskrit for massages, then the Sanskrit for a secluded place.
\end{quote}

\subsection*{You'll want to know: \sk{स्मयते}}
\textbf{\sk{स्मयते}} — \emph{smayate} — \textquotedblleft{}smiles\textquotedblright{}.

\begin{quote}
\sk{शिशुः} \sk{मातरं} \sk{पश्यति}, \sk{स्मयते} \sk{च।}
\end{quote}

\begin{quote}
\emph{śiśuḥ mātaraṃ paśyati, smayate ca.}
\end{quote}

\begin{quote}
The baby sees its mother and smiles.
\end{quote}

Past (she …): \sk{सा} \sk{स्मितवती}, \emph{sā smitavatī}.

\begin{cousinweb}
\textbf{\sk{स्मयते}} (\emph{smayate}) comes from the root √smi, \textquotedblleft{}to smile\textquotedblright{}, the root of \textbf{\sk{विस्मयः}} (MW).
\end{cousinweb}

\begin{grammarlens}[title={What you've built}]
One more word for describing what you lived and felt.
\end{grammarlens}

\subsection*{Your turn}
\begin{itemize}
\raggedright
  \item \emph{Say it:} \emph{smayate}
  \item \emph{Say it:} \emph{smayate}, once more
  \item \emph{Say it:} \emph{śiśuḥ mātaraṃ paśyati, smayate ca.} — The baby sees its mother and smiles.
  \item \emph{Recall:} say the Sanskrit for massages, then the Sanskrit for a secluded place, then say \emph{smayate} again
\end{itemize}

\begin{tcolorbox}[breakable,colback=teal!4,colframe=teal!35!black,title={Before you move on}]
What does \sk{स्मयते} mean? (\textquotedblleft{}Smiles\textquotedblright{}.) Say it once more.
\end{tcolorbox}
\section[\texorpdfstring{\sk{उपहसति} (upahasati)}{उपहसति (upahasati)}]{\sk{उपहसति} (upahasati) — laughs at}

\label{lesson:SA-C457-upahasati}



\begin{quote}
Before the new one: say the Sanskrit for a secluded place, then the Sanskrit for smiles.
\end{quote}

\subsection*{You'll want to know: \sk{उपहसति}}
\textbf{\sk{उपहसति}} — \emph{upahasati} — \textquotedblleft{}laughs at\textquotedblright{}.

\begin{quote}
\sk{बालाः} \sk{मूर्खं} \sk{वानरम्} \sk{उपहसन्ति।}
\end{quote}

\begin{quote}
\emph{bālāḥ mūrkhaṃ vānaram upahasanti.}
\end{quote}

\begin{quote}
The children laugh at the silly monkey.
\end{quote}

Past (he …): \sk{सः} \sk{मित्रम्} \sk{उपहसितवान्}, \emph{saḥ mitram upahasitavān}.

\begin{cousinweb}
\textbf{\sk{उपहसति}} (\emph{upahasati}) is built from \emph{upa}, \textquotedblleft{}towards\textquotedblright{}, and the root √has, \textquotedblleft{}to laugh\textquotedblright{}, the root of \textbf{\sk{हसति}}; MW gives \textquotedblleft{}to laugh at, deride, ridicule\textquotedblright{} (MW).
\end{cousinweb}

\begin{grammarlens}[title={What you've built}]
One more word for describing what you lived and felt.
\end{grammarlens}

\subsection*{Your turn}
\begin{itemize}
\raggedright
  \item \emph{Say it:} \emph{upahasati}
  \item \emph{Say it:} \emph{upahasati}, once more
  \item \emph{Say it:} \emph{bālāḥ mūrkhaṃ vānaram upahasanti.} — The children laugh at the silly monkey.
  \item \emph{Recall:} say the Sanskrit for a secluded place, then the Sanskrit for smiles, then say \emph{upahasati} again
\end{itemize}

\begin{tcolorbox}[breakable,colback=teal!4,colframe=teal!35!black,title={Before you move on}]
What does \sk{उपहसति} mean? (\textquotedblleft{}Laughs at\textquotedblright{}.) Say it once more.
\end{tcolorbox}
\section[\texorpdfstring{\sk{परिहासः} (parihāsaḥ)}{परिहासः (parihāsaḥ)}]{\sk{परिहासः} (parihāsaḥ) — a jest}

\label{lesson:SA-C457-parihasah}



\begin{quote}
Before the new one: say the Sanskrit for smiles, then the Sanskrit for laughs at.
\end{quote}

\subsection*{You'll want to know: \sk{परिहासः}}
\textbf{\sk{परिहासः}} — \emph{parihāsaḥ} — \textquotedblleft{}a jest\textquotedblright{}.

\begin{quote}
\sk{अरुणस्य} \sk{परिहासेन} \sk{सर्वे} \sk{हसन्ति।}
\end{quote}

\begin{quote}
\emph{aruṇasya parihāsena sarve hasanti.}
\end{quote}

\begin{quote}
Everyone laughs at Arun's joke.
\end{quote}

A masculine noun.

\begin{cousinweb}
\textbf{\sk{परिहासः}} (\emph{parihāsaḥ}) is built from \emph{pari}, \textquotedblleft{}around\textquotedblright{}, and the root √has, \textquotedblleft{}to laugh\textquotedblright{}; MW gives \textquotedblleft{}jesting, joking\textquotedblright{} and \textquotedblleft{}a jest, joke\textquotedblright{} (MW).
\end{cousinweb}

\begin{grammarlens}[title={What you've built}]
One more word for describing what you lived and felt.
\end{grammarlens}

\subsection*{Your turn}
\begin{itemize}
\raggedright
  \item \emph{Say it:} \emph{parihāsaḥ}
  \item \emph{Say it:} \emph{parihāsaḥ}, once more
  \item \emph{Say it:} \emph{aruṇasya parihāsena sarve hasanti.} — Everyone laughs at Arun's joke.
  \item \emph{Recall:} say the Sanskrit for smiles, then the Sanskrit for laughs at, then say \emph{parihāsaḥ} again
\end{itemize}

\begin{tcolorbox}[breakable,colback=teal!4,colframe=teal!35!black,title={Before you move on}]
What does \sk{परिहासः} mean? (\textquotedblleft{}A jest\textquotedblright{}.) Say it once more.
\end{tcolorbox}
\section[\texorpdfstring{\sk{अट्टहासः} (aṭṭahāsaḥ)}{अट्टहासः (aṭṭahāsaḥ)}]{\sk{अट्टहासः} (aṭṭahāsaḥ) — a loud laugh}

\label{lesson:SA-C457-attahasah}



\begin{quote}
Before the new one: say the Sanskrit for laughs at, then the Sanskrit for a jest.
\end{quote}

\subsection*{You'll want to know: \sk{अट्टहासः}}
\textbf{\sk{अट्टहासः}} — \emph{aṭṭahāsaḥ} — \textquotedblleft{}a loud laugh\textquotedblright{}.

\begin{quote}
\sk{पिता} \sk{अट्टहासं} \sk{करोति}, \sk{बालाः} \sk{अपि} \sk{हसन्ति।}
\end{quote}

\begin{quote}
\emph{pitā aṭṭahāsaṃ karoti, bālāḥ api hasanti.}
\end{quote}

\begin{quote}
Father lets out a loud laugh, and the children laugh too.
\end{quote}

A masculine noun.

\begin{cousinweb}
\textbf{\sk{अट्टहासः}} (\emph{aṭṭahāsaḥ}) is built from \emph{aṭṭa}, \textquotedblleft{}high, lofty\textquotedblright{}, and \emph{hāsa}, \textquotedblleft{}laughter\textquotedblright{}, from the root √has, \textquotedblleft{}to laugh\textquotedblright{} (MW).
\end{cousinweb}

\begin{grammarlens}[title={What you've built}]
One more word for describing what you lived and felt.
\end{grammarlens}

\subsection*{Your turn}
\begin{itemize}
\raggedright
  \item \emph{Say it:} \emph{aṭṭahāsaḥ}
  \item \emph{Say it:} \emph{aṭṭahāsaḥ}, once more
  \item \emph{Say it:} \emph{pitā aṭṭahāsaṃ karoti, bālāḥ api hasanti.} — Father lets out a loud laugh, and the children laugh too.
  \item \emph{Recall:} say the Sanskrit for laughs at, then the Sanskrit for a jest, then say \emph{aṭṭahāsaḥ} again
\end{itemize}

\begin{tcolorbox}[breakable,colback=teal!4,colframe=teal!35!black,title={Before you move on}]
What does \sk{अट्टहासः} mean? (\textquotedblleft{}A loud laugh\textquotedblright{}.) Say it once more.
\end{tcolorbox}
\section[\texorpdfstring{\sk{विनोदः} (vinodaḥ)}{विनोदः (vinodaḥ)}]{\sk{विनोदः} (vinodaḥ) — amusement}

\label{lesson:SA-C457-vinodah}



\begin{quote}
Before the new one: say the Sanskrit for a jest, then the Sanskrit for a loud laugh.
\end{quote}

\subsection*{You'll want to know: \sk{विनोदः}}
\textbf{\sk{विनोदः}} — \emph{vinodaḥ} — \textquotedblleft{}amusement\textquotedblright{}.

\begin{quote}
\sk{विनोदाय} \sk{बालाः} \sk{उद्याने} \sk{क्रीडन्ति।}
\end{quote}

\begin{quote}
\emph{vinodāya bālāḥ udyāne krīḍanti.}
\end{quote}

\begin{quote}
The children play in the garden for fun.
\end{quote}

A masculine noun.

\begin{cousinweb}
\textbf{\sk{विनोदः}} (\emph{vinodaḥ}) is built from \emph{vi} and the root √nud, \textquotedblleft{}to push\textquotedblright{}, the root of \textbf{\sk{नुदति}}: driving away dullness; MW gives \textquotedblleft{}removal\textquotedblright{} first and then \textquotedblleft{}diversion, pastime, pleasure\textquotedblright{} (MW).
\end{cousinweb}

\begin{grammarlens}[title={What you've built}]
One more word for describing what you lived and felt.
\end{grammarlens}

\subsection*{Your turn}
\begin{itemize}
\raggedright
  \item \emph{Say it:} \emph{vinodaḥ}
  \item \emph{Say it:} \emph{vinodaḥ}, once more
  \item \emph{Say it:} \emph{vinodāya bālāḥ udyāne krīḍanti.} — The children play in the garden for fun.
  \item \emph{Recall:} say the Sanskrit for a jest, then the Sanskrit for a loud laugh, then say \emph{vinodaḥ} again
\end{itemize}

\begin{tcolorbox}[breakable,colback=teal!4,colframe=teal!35!black,title={Before you move on}]
What does \sk{विनोदः} mean? (\textquotedblleft{}Amusement\textquotedblright{}.) Say it once more.
\end{tcolorbox}
